\subsection{All higher-precision relative bilinear forms}

\begin{statement}{A.129}{Proved}
{The semi-invariant operator module and its exact coefficient maps}
Let \(a\in\Q^\times\) satisfy \(v_3(a)=1\), let \(n,r\ge1\),
and retain the actual relative norm torus \(T_n\) of A.115.
Put
\[
 N=3^n,\quad d=N/3,\quad R_r=\Z/3^r\Z,\quad
 C_r=T_n[3^r],\quad m=\max(n,r).
 \tag{A.129.1}
\]
Use compatible roots of unity to write the actual splitting group as
\(G_{n,r}=C_N\rtimes U_m\), with
\(U_m=(\Z/3^m\Z)^\times\).
Let \(\rho_r:C_r\to C_1\) be multiplication by \(3^{r-1}\),
and let \(\iota_r:\mu_3\hookrightarrow\mu_{3^r}\) be inclusion.
Every equivariant bilinear form is alternating, and the map
\[
 \begin{split}
 \operatorname{Hom}_{G_\Q}(\textstyle\bigwedge^2 C_1,\mu_3)
 &\xrightarrow{\ \sim\ }
       \operatorname{Hom}_{G_\Q}(C_r\otimes_{R_r}C_r,\mu_{3^r}),\\
 B&\longmapsto\iota_r\circ B\circ(\rho_r,\rho_r)
 \end{split}
 \tag{A.129.2}
\]
is an isomorphism of abelian groups.  In particular the group on the
right is \((\Z/3\Z)^n\), and every form has values in the literal
order-three subgroup of \(\mu_{3^r}\).

Here is the integral operator calculation giving the map.
Put
\[
 \mathcal O_n=\Z_3[x]/(\Phi_N(x)),\qquad
 \Phi_N(x)=1+x^d+x^{2d},\qquad
 \sigma_u(x)=x^u.
 \tag{A.129.3}
\]
Let \(\omega(u)\in\{1,-1\}\) lift \(u\bmod3\), and let
\(\chi_r(u)=u\bmod3^r\).
Under the perfect underlying block pairing of A.115, the group
of forms in (A.129.2) is exactly
\[
 \begin{split}
 \{P\in\mathcal O_n/3^r:\ 
        \sigma_u(P)=\omega(u)\chi_r(u)P\text{ for all }u\in U_m\}
 &=\ker(\sigma_2+2:\mathcal O_n/3^r\to\mathcal O_n/3^r)\\
 &=3^{r-1}\ker(\sigma_2-1:\mathcal O_n/3\to\mathcal O_n/3).
 \end{split}
 \tag{A.129.4}
\]
The last notation means the image of the kernel under
\(\mathcal O_n/3\to\mathcal O_n/3^r\), \(P\mapsto3^{r-1}\widetilde P\).
It is independent of the chosen lift.

The coefficient transitions retain their literal maps.
For \(r\ge2\), reducing the value of any such form by
\(\mu_{3^r}\xrightarrow{z\mapsto z^3}\mu_{3^{r-1}}\)
gives zero.  Restricting its two arguments to the subgroup
\(C_{r-1}\hookrightarrow C_r\) also gives zero.
For an actual continuous class \(\gamma\in H^1(\Q_v,C_r)\),
the associated quadratic class is
\[
 \tfrac12\bigl(\iota_r B(\rho_r,\rho_r)\bigr)_*
                         (\gamma\cup\gamma)
  =(\iota_r)_*\!\left(
       \tfrac12 B_*((\rho_r)_*\gamma\cup(\rho_r)_*\gamma)\right).
 \tag{A.129.5}
\]
Consequently these higher-precision bilinear reciprocity tests
are exactly the mod-three tests under the displayed maps.
They do not characterize a higher-precision local cyclic image:
an explicit relative norm-one class at precision nine below
has every displayed bilinear value zero and has noncyclic
local image.
\end{statement}
\begin{proof}
\emph{The actual module and its pairing character.}
The geometric coordinate norm gives
\[
 C_r=I_r\otimes\mu_{3^r},\qquad
 I_r=\ker\bigl(R_r[C_N]\longrightarrow R_r[C_d]\bigr).
 \tag{A.129.6}
\]
The map in (A.129.6) reduces root labels modulo \(d\);
equivalently it sums the three coordinates in each block.
Multiplication by \(x^d-1\) identifies
\(\mathcal O_n/3^r\) with \(I_r\): its image is the kernel,
its annihilator in \(R_r[C_N]\) is generated by \(\Phi_N\),
and the basis \((x^d-1)x^j\), \(0\le j<2d\), is a basis
of the kernel.  Translation is multiplication by \(x\).
An endomorphism commuting with this cyclic action is therefore
multiplication by one unique \(P\in\mathcal O_n/3^r\).
Arithmetic conjugation on these multiplication operators is
\(P(x)\mapsto P(x^u)\); the factors by which the generator
\(x^d-1\) and the Tate twist change cancel in this conjugation.

On the two independent coordinates in each block put
\[
       B_*(v,w)=\sum_{h=0}^{d-1}
                         (v_hw_{h+d}-v_{h+d}w_h).
 \tag{A.129.7}
\]
This is a perfect alternating pairing of underlying free
\(R_r\)-modules.  Translation permutes blocks and rotates
their three coordinates, with determinant one.
The permutation of a block induced by a unit \(u\) has
determinant \(\omega(u)\).  The two Tate twists contribute
\(\chi_r(u)^2\).  Thus the exact arithmetic identity is
\[
                B_*(uv,uw)=\chi_r(u)^2\omega(u)B_*(v,w).
 \tag{A.129.8}
\]
For \(r>1\) this pairing character must be retained.
For example when \(n=1,r=2\), the unit \(u=4\) fixes
every embedding and multiplies \(C_r\) by four.
A form with target \(\mu_9\) is forced to satisfy
\((16-4)B=0\), hence \(3B=0\).

Every translation-equivariant bilinear form has the unique
expression \(B_P(v,w)=B_*(v,Pw)\).
By (A.129.8), its target character is \(\chi_r\) exactly when
\[
 u^{-1}Pu=\chi_r(u)^{-1}\omega(u)P,
       \quad\text{or equivalently}\quad
        \sigma_u(P)=\chi_r(u)\omega(u)P.
 \tag{A.129.9}
\]
The unit two generates \(U_m\), and \(\omega(2)=-1\),
so this is exactly \((\sigma_2+2)P=0\).
This proves the first equality in (A.129.4).

\emph{The complete integral Smith calculation.}
Write \(D=\sigma_2+2\) on the free \(\Z_3\)-module
\(\mathcal O_n\) of rank \(2d\).
The characteristic polynomial of \(\sigma_2\) over
\(\Q_3\) is \(X^{2d}-1\).
One explicit verification uses the separable embedding map
\[
 \Q_3(\zeta_N)\otimes_{\Q_3}\overline{\Q}_3
       \longrightarrow
        \prod_{u\in U_n}\overline{\Q}_3,
             \qquad b\otimes c\longmapsto(\sigma_u(b)c)_u.
 \tag{A.129.10}
\]
It is an isomorphism, and \(\sigma_2\) permutes its factors
in one cycle of length \(2d\).
Equally, the ordinary permutation modules have the exact
sequence
\(0\to\Z_3[C_d]\xrightarrow{\Phi_N}\Z_3[C_N]
  \to\mathcal O_n\to0\).
The orbit of labels of valuation zero is the one additional
cycle, giving the same characteristic polynomial.
It follows that
\[
            \det D=2^{2d}-1,\qquad
                      v_3(\det D)=n.
 \tag{A.129.11}
\]
The valuation identity follows from
\(v_3(4^{3^{n-1}}-1)=n\).

Modulo three the operator \(D\) is \(\sigma_2-1\).
Here \(\mathcal O_n/3=\F_3[Z]/Z^{2d}\), with \(x=1+Z\).
A.115 computes its full invariant space, of dimension \(n\),
with basis
\[
        z^{d-3^j},\quad 0\le j<n,\qquad
                            z=Z^2/(1+Z).
 \tag{A.129.12}
\]
There are therefore exactly \(n\) nonunit Smith entries
of \(D\) over \(\Z_3\).  Their valuations are positive
integers whose sum is \(n\), by (A.129.11).
Each has valuation exactly one; every other Smith entry
is a unit.  In Smith coordinates, the kernel modulo
\(3^r\) is consequently killed by three.
Returning by the invertible integral change of basis
shows that every element of this kernel is divisible
by \(3^{r-1}\) in \(\mathcal O_n/3^r\).
Writing it as \(3^{r-1}\widetilde P\) gives
\(D\widetilde P=0\bmod3\), proving the last equality
of (A.129.4).

Formula (A.129.12) and the actual identification (A.129.7)
now give every form as
\(3^{r-1}B_*(v,\widetilde P w)\).
It depends only on the reductions of its arguments modulo
three, and its value is the inclusion of the corresponding
form of A.115 into \(\mu_{3^r}\).
All those mod-three forms are alternating; the lifted
forms are therefore alternating as well.
This proves (A.129.2), including its assertion about all
bilinear forms, not only those initially assumed alternating.

\emph{Transitions and the cohomology comparison.}
In compatible additive root coordinates, \(\rho_r\)
is reduction modulo three and \(\iota_r\) is multiplication
by \(3^{r-1}\).  The value transition to level \(r-1\)
reduces that multiple modulo \(3^{r-1}\), giving zero.
The inclusion of \(C_{r-1}\) in \(C_r\) multiplies
each coordinate by three; hence \(\rho_r\) is zero
on that subgroup.  These are the two asserted zero maps.
Naturality of the continuous cup product gives (A.129.5)
at the cochain level and therefore in cohomology.
The same identities show that an integral equivariant
bilinear pairing on \(T_3T_n\) with values in \(\Z_3(1)\)
is zero: all its compatible finite reductions vanish
because every coefficient transition just computed is zero.

\emph{An actual local precision-nine comparison.}
Take \(n=1\), \(r=2\), and \(q=321193\).
The exact splitting and primitive-unit certificate in A.106
shows that \(q\) splits \(\Q(3^{1/3},\mu_9)\) and
that \(g=5\) generates the residue multiplicative group.
Thus \(T_1(\Q_q)\) is the literal product-one subgroup
of \((\Q_q^\times)^3\).
The actual element
\[
                       (q,\ q^{-1}g^3,\ g^{-3})
 \tag{A.129.13}
\]
has product one.
Its tame and unramified Kummer vectors modulo nine are
\(t=(1,-1,0)\) and \(c=(0,3,-3)\).
Their mod-three reductions are \(\overline t\) and zero,
so all forms (A.129.2) evaluate to zero.
Nevertheless its local homomorphism has noncyclic image.
Indeed \(t\) has order nine, so a cyclic subgroup containing
it must be the group it generates, as the whole coefficient
module has exponent nine.  But \(c\ne zt\) for every
\(z\in\Z/9\): the first coordinate forces \(z=0\),
contradicting the second coordinate.
This supplies the promised actual comparison without
identifying vanishing bilinear tests with a higher cyclic
image condition.
\end{proof}

The retained program
\path{code/relative_precision_forms_certificate.py}
computes the integral substitution matrices and their Smith
entries through \(n=5\), checks thirty exact kernel sizes
through precision \(3^6\), and verifies the literal block
actions, all displayed form generators and forty-eight
coefficient/source transition identities.
Its saved output also checks the local precision-nine
noncyclic example.  The proof above gives the assertions
at every \(n\) and every \(r\).
